\label{exp:cec01}

\subsubsection*{Постановка}
\textbf{Физический вопрос.}
На микроуровне~$\Mlayer$ локальный закон~$g$ строго обратим ($A_{13}$) и симметричен относительно обращения знака фазы, однако наш каузальный пузырь на макроуровне~$\Tlayer$ содержит избыток барионов ($n=+1$) над антибарионами ($n=-1$).
Способен ли детерминированный шаг~$g$ на трёхмерной решётке~\latGCKshort\ без посаженных начальных вихрей ($b(\tau=0)\equiv 0$) породить устойчивые барионные солитоны ($b=1$) из 9-битного начального вектора зародыша $\mathrm{IV}_{\mathrm{bubble}} = 2^{\vevHierarchySteps} + d_{\mathrm{soft}} = 263 = \texttt{0b1\_00000111}$ (\cref{thm:cosmo-9bit-iv-baryon}), и меняет ли знак рождающегося мира инверсия старшего бита хиральности ($\mathrm{IV} = 7 = \texttt{0b0\_00000111}$)?

\textbf{Что проверяем.}
\begin{enumerate}[label=\textbf{(\alph*)},leftmargin=*,itemsep=0.25em]
\item Прямой прогон 3D-решётки~\latGCKshort\ при $\mathrm{IV} = 263$ ($b_{\mathrm{chir}}=1$, $\mathrm{mask}_8 = 7$): нуклеация вихревых ядер ($b=1$, \texttt{born=1}) и положительный суммарный топологический заряд $\sum n > 0$.
\item Зеркальный контроль $\mathrm{IV} = 7$ ($b_{\mathrm{chir}}=0$, $\mathrm{mask}_8 = 7$): инверсия знака заряда ($\sum n < 0$, антиматерия).
\item Нулевые контроли: однородный сдвиг ($\mathrm{mask}_8 = 0$) и плоская одномерная стенка (\texttt{wall}, $\mathrm{mask}_8 = 15$), где контурный интеграл фазы равен нулю ($\oint\mathrm{d}\varphi = 0 \implies \sum b \equiv 0$).
\end{enumerate}

\textbf{Рабочая гипотеза.}
Только трёхгранный узел Киббла--Зурека на стыке трёх ортогональных октантов барионного треугольника~$K_3$ ($\mathrm{mask}_8 = 7 = d_{\mathrm{soft}}$) создаёт ненулевую голономию и переводит фазовый импульс через порог дискретности $\Delta\varphi_{\mathrm{disc}} = 41$, а знак $\sum n$ однозначно определяется старшим битом~$b_{\mathrm{chir}}$.

\textbf{Наблюдаемые.}
Число родившихся ядер $N_{\mathrm{born}} = \sum_x b(x)$, суммарный топологический заряд $Q_{\mathrm{tot}} = \sum_x n(x)$ и булев индикатор нуклеации \texttt{born} на такте $\tau = 32$ (табл.~\ref{tab:exp-ce-c-01}).

\textbf{Критерий согласия с теорией.}
Детерминированное выполнение $\texttt{born}=1$, $Q_{\mathrm{tot}} > 0$ при $\mathrm{IV}=263$; $Q_{\mathrm{tot}} < 0$ при $\mathrm{IV}=7$; строгое тождество $N_{\mathrm{born}} \equiv 0$ при $\mathrm{mask}_8 \in \{0, 15\}$ (тип~B, целочисленная арифметика).

\subsubsection*{Теоретическая часть}
По теореме~\ref{thm:cosmo-9bit-iv-baryon} три координатные плоскости вокруг центра нуклеации делят пространство на $2^d = 8 = \vevHierarchySteps$ октантов, что в сумме с битом хиральности спинора $b_{\mathrm{chir}}\in\{0,1\}$ исчерпывает весь 9-битный регистр ячейки Бекештейна ($\dim_2(\mathrm{IV}) = 1 + 8 = 9 = \lfloor\caBitBudget\rfloor$).
Плоская стенка меняет фазу лишь вдоль одной оси и даёт $\oint_{\partial f}\Delta\varphi = 0$ на каждой треугольной плакетке~$K_3$, тогда как трёхгранный угол $\mathrm{mask}_8 = 7$ концентрирует скачок фазы на общей вершине трёх плоскостей $\{XY, YZ, XZ\}$.

\textbf{Проверка размерностей.}
Все входы и выходы опыта строго безразмерны:
\[
  [\mathrm{IV}] = 1\quad(\mathbb{Z}_{512}),\qquad
  [\tau/\htick] = 1,\qquad
  [b(x)] = 1\quad(\{0,1\}),\qquad
  \left[\sum_x n(x)\right] = 1\quad(\mathbb{Z}),
\]
а перевод к кулонам в СИ осуществляется умножением целого~$n$ на элементарный заряд~$e_0$ ($[Q_{\mathrm{SI}}] = [n\,e_0] = \mathrm{I\cdot T}$).

\subsubsection*{Ход эксперимента}
Скрипт воспроизведения: \texttt{python scripts/run\_phase\_iv\_to\_i.py --fcc}
(функция генерации сида в \texttt{mt\_ca/seeds.py}).

\begin{enumerate}[label=\textbf{Шаг \arabic*:},leftmargin=*,itemsep=0.35em]
\item Инициализировать 3D-решётку~\latGCKshort\ ($16^3$) в режиме кипящего пола~$A_5$ с $b(\tau=0)\equiv 0$.
\item Подать 9-битный вектор $\mathrm{IV}\in\{263, 7, 0, 15\}$ и выполнить $32$ обратимых шага~$g$.
\item Снять $N_{\mathrm{born}} = \sum_x b(x)$ и $Q_{\mathrm{tot}} = \sum_x n(x)$ в табл.~\ref{tab:exp-ce-c-01}.
\end{enumerate}

\begin{table}[htbp]
\centering
\footnotesize
\setlength{\tabcolsep}{4pt}
\caption{Нуклеация барионной асимметрии из 9-битного вектора $\mathrm{IV}\in\mathbb{Z}_{512}$ на 3D~\latGCKshort\ (CE-C-01)}
\label{tab:exp-ce-c-01}
\begin{tabular}{@{}llllr@{}}
\toprule
Режим сида & $\mathrm{IV}\in\mathbb{Z}_{512}$ & $b_{\mathrm{chir}}$ & Маска октантов $\mathrm{mask}_8$ & Итог к $\tau=32$ ($\texttt{born}$, $\sum n$) \\
\midrule
Наш пузырь ($K_3$, материя) & $263$ (\texttt{0b1\_00000111}) & $1$ & $7 = d_{\mathrm{soft}}$ ($3$ октанта $K_3$) & $\texttt{born}=1,\ \sum n = +2 > 0$ \\
Зеркальный (антиматерия) & $7$ (\texttt{0b0\_00000111}) & $0$ & $7 = d_{\mathrm{soft}}$ ($3$ октанта $K_3$) & $\texttt{born}=1,\ \sum n = -2 < 0$ \\
Однородный лёд (контроль) & $0$ (\texttt{0b0\_00000000}) & $0$ & $0$ (без градиента) & $\texttt{born}=0,\ \sum b \equiv 0$ \\
Плоская стенка (\texttt{wall}) & $15$ (\texttt{0b0\_00001111}) & $0$ & полупространство $x\ge c_x$ & $\texttt{born}=0,\ \sum b \equiv 0$ \\
\bottomrule
\end{tabular}
\end{table}

\subsubsection*{Анализ, расчёт погрешности и выводы}
Поскольку ядро работает в точной модулярной арифметике $\mathbb{Z}_{512}^4$, статистическая погрешность типа~A равна нулю ($u_A = 0$), а погрешность типа~B отсутствует на целочисленных инвариантах ($b\in\{0,1\}$, $n\in\mathbb{Z}$).
Таблица~\ref{tab:exp-ce-c-01} доказывает, что трёхгранный узел $\mathrm{IV}=263$ детерминированно рождает материю ($\sum n = +2$), а замена старшего бита $\mathrm{IV}=7$ зеркально рождает антиматерию ($\sum n = -2$).

\textbf{Вывод.}
Барионная асимметрия нашего пузыря однозначно кодируется 9-битным сидом $\mathrm{IV}=263\in\mathbb{Z}_{512}$ без нарушения обратимости закона~$g$.
\textbf{Статус записи:} опыт закрыт по табл.~\ref{tab:exp-ce-c-01}.
